% Transcription — fonds Alexandre Grothendieck, shelfmark 89, batch 1 (pages 1-20).
% Established from the facsimile archives/batches/89/batch-01.pdf (a mirror of
% https://grothendieck.umontpellier.fr/89.pdf), per the skill
% transcribe-grothendieck. The folder is twenty pages and this is its only
% batch.
%
% Nineteen of the twenty leaves are transcribed. Page 1 is a cover carrying
% only « les ensembles 3-3 » in pencil; it takes no \page marker and its words
% become the section heading. The inventory brackets both halves of its title,
% but « Ensembles 3-3 » is his (the cover) and « hexagone combinatoire » is his
% too (page 2, first line). There is no pagination of his own; the runs are
% followed by his roman numerals I (page 2) and II (page 9) and, for the last
% two leaves, by their subject.
%
% Four runs. Pages 2-8 (« I »): a combinatorial hexagon H = (S, A, R) amounts
% to an orientation omega_H, a set Delta of three diagonals and a pair T of
% inscribed triangles, with S = Delta x T; the groupoid of combinatorial
% hexagons is equivalent to Ens_3 x Ens_2 (his margin). He then studies a
% six-element set with a (3,3)-partition and a transitive system of isos
% between the orientations omega_{S_i}: its automorphism group G has order 36;
% a subordinate hexagonal structure is an « antipodisme » sigma; the three
% antipodismes generate G' ≅ S_I, presented by sigma^2 = sigma'^2 =
% (sigma sigma')^3 = 1; G' acts trivially on the Gamma_omega-torsor T = wedge
% S_i, so G/G' (order 6) acts on it. Pages 7-8 restart with the full
% automorphism group (order 72) of a (3,3)-partitioned set, whose involutive
% splittings S' form a set of six elements surjecting 3-to-1 onto T' = wedge
% omega_{S_i}; a hexagon is a pair (S, R) with a point of S'. The run breaks
% off mid-sentence. Pages 9-14 (« II »): a plane V over F_3 has four lines;
% a pair T of them and the complementary pair T' give two decompositions of V
% and of any V-torsor P (nine points) as a product of two three-point sets —
% two pairs of « orthogonal » (3,3,3)-partitions; hence an equivalence (8)
% between (3,3)-partitioned six-element sets with bijections omega_{P_i} ≅
% D_i^* indexed by T and by T', quasi-inverse to itself under T <-> T';
% combinatorial squares Q correspond to planes over F_3 with a pair of lines
% (9)-(10), the complementary pair gives the dual square DQ, and alpha_Q
% alpha_{DQ} = -id = V(a_Q), the antipodism of Q (13). Pages 15-18: for V
% free of rank 2 over a ring Lambda, the aim is to recover V from P(V)^! with
% the torsors T_x = D_x^*, omega_V = det(V)^*, the disjointness relation and
% the isomorphisms T_x wedge T_y ≅ omega_V; for u, v, w pairwise disjoint,
% u + v + w = 0 iff u wedge v = v wedge w = w wedge u, which gives addition on
% disjoint pairs; the statement is for a field, and the list of properties
% needed breaks off after a). Pages 19-20, a different run in a faster hand:
% combinatorial polygons from three points of view (ordered set S u A, edges
% A in P_2(S), a pair of opposite circular orders), orientations of edges and
% contours, and a theorem: a polygon has exactly two orientations, each
% determining a unique bijection f_omega sending s along its outgoing edge.
%
% Reading hazards fixed once here. « syst. trans. d'isom. » is his; « ens. »,
% « iso », « autom. », « corr. », « quelc. » are abbreviations kept as
% written. He writes ∧ for the contracted product of torsors, sometimes with
% the group underneath (wedge_{±1}); it is kept. His underlining is set in
% bold. The circled tall S of the symmetric group is set \mathfrak{S}; the
% script G of page 7 is \mathcal{G}; « Tors » is set \mathbf{Tors}. Page 11's
% double-headed ≈ strokes needed a `leftrightarrow` arrow option, added to
% scripts/render.mjs in the same commit. Pages 4 (a fan-shaped marginal note
% at some sixty degrees), 17 (two slanted marginal notes) and 20 (a slanted
% margin with a long elbowed connector, three cancelled lines) are the worst
% of the hand; what they will not support is left \ill{}.
%
% Readings flagged and not settled: « structure unisignée » opening page 4;
% « inverseurs » on page 4; the reading of the term closing page 5; « plan »
% on page 9 (struck, with nothing substituted); « carré » on page 13, a
% underlined word that could pass for « couvé », settled by the diagonals;
% the numeral of (12) on page 14, written as (11); « torseur » on page 15;
% « automorphismes » in c) on page 20.
%
% Pass: Opus 5 (claude-opus-5), 2026-09-13 — first pass, unchecked against the pages by a human.
\documentclass[11pt,a4paper]{article}
% Le préambule commun à toutes les transcriptions du fonds Grothendieck.
%
% Il tient en une idée : **l'incertitude se compose, elle ne se résout pas.**
% Une transcription qui lisse un mot douteux a détruit la seule chose pour
% laquelle on la faisait — un lecteur doit pouvoir savoir, à chaque endroit,
% ce qui était sur le feuillet et ce qui a été ajouté. Les six macros qui
% suivent sont donc l'appareil critique, et non de la décoration.
%
% Les mêmes commandes sont reconnues par `scripts/render.mjs`, qui produit la
% vue de lecture du site. Une macro ajoutée ici doit l'être aussi là-bas,
% faute de quoi le rendu échouera bruyamment — ce qui est voulu : un rendu
% silencieusement faux est pire qu'un rendu absent.

\ProvidesPackage{grothendieck}[2026/08/08 Transcriptions du fonds Grothendieck]

\usepackage{amsmath,amssymb,amsthm}
\usepackage{mathrsfs}
\usepackage{stmaryrd}
% Les diagrammes commutatifs sont partout dans ce fonds : c'est la langue
% même de la géométrie algébrique de Grothendieck, et une transcription qui
% les rendrait en prose ne transcrirait rien.
\usepackage{tikz-cd}
% Français par défaut — c'est la langue des feuillets — mais l'anglais est
% chargé aussi : la traduction et le résumé sont des documents anglais, et la
% typographie française (espaces avant les deux-points) y serait une faute.
% Ces documents basculent par \selectlanguage{english} en tête de corps.
\usepackage[english,french]{babel}
\usepackage{fontspec}
\usepackage{geometry}
\geometry{a4paper,margin=25mm,marginparwidth=28mm}
\usepackage{marginnote}
\usepackage{xcolor}
% ulem, not soul. `soul`'s \st and \ul are built on a character-by-character
% scan that XeLaTeX's font machinery breaks: the document fails with an
% undefined control sequence rather than degrading. ulem does the same two jobs
% and is Unicode-engine safe. `normalem` stops it hijacking \emph, which the
% transcriptions use for Grothendieck's own emphasis.
\usepackage[normalem]{ulem}
\usepackage{hyperref}
\hypersetup{colorlinks=true,linkcolor=black,urlcolor=black,pdfborder={0 0 0}}

% --- Métadonnées du lot -----------------------------------------------------
% Reprises telles quelles par le script de rendu, qui les affiche en tête de
% la vue de lecture. Elles fixent ce que le fichier prétend couvrir : sans
% elles, une transcription flotte sans point d'attache dans un fonds de seize
% mille feuillets.

\newcommand{\folder}[1]{\gdef\grothFolder{#1}}
\newcommand{\batch}[1]{\gdef\grothBatch{#1}}
\newcommand{\leaves}[2]{\gdef\grothFirst{#1}\gdef\grothLast{#2}}
\newcommand{\foldertitle}[1]{\gdef\grothFoldertitle{#1}}
\gdef\grothFolder{?}\gdef\grothBatch{?}\gdef\grothFirst{?}\gdef\grothLast{?}\gdef\grothFoldertitle{}

% --- Le résumé --------------------------------------------------------------
%
% Il ouvre la lecture modernisée, sous le titre « Résumé », et il est composé
% en petit. Ce n'est pas un ornement : c'est le seul passage du document qui
% ne soit pas des mathématiques, et un lecteur doit pouvoir le reconnaître
% avant de l'avoir lu — puis le sauter s'il connaît déjà le sujet, ou s'y
% arrêter s'il ne le connaît pas. Le filet qui le suit marque la reprise du
% corps.

\newenvironment{resume}
  {\small\setlength{\parskip}{0.45em}}
  {\par\medskip\hrule\medskip}

% --- Mots-clés ---------------------------------------------------------------
%
% En anglais, en clôture du résumé de la lecture modernisée : le vocabulaire
% moderne sous lequel chercher ce que les feuillets construisent. Le manifeste
% du site (`npm run manifest`) les extrait pour étiqueter la cote entière —
% cette ligne est la source unique des tags, il n'y a pas de fichier à côté.
\newcommand{\keywords}[1]{\par\smallskip\noindent
  {\footnotesize\textsc{Keywords} --- \textit{#1}}\par}

% --- L'appareil critique ----------------------------------------------------

% Le numéro de feuillet, en marge. C'est l'ancre qui permet de régler un
% désaccord sur une formule en regardant *un* feuillet plutôt qu'une chemise
% de six cents. Le site s'en sert aussi pour tourner les pages du fac-similé
% quand on fait défiler la transcription.
\newcommand{\leaf}[1]{%
  \marginnote{\footnotesize\color{gray!70}\textsf{#1}}%
  \label{leaf:#1}\ignorespaces}

% Illisible : on ne devine pas. Un mot inventé qui se lit comme les autres est
% le pire résultat possible.
\newcommand{\ill}{\textcolor{red!60!black}{[\,\ldots\,]}}

% Lecture douteuse : le mot est donné, et signalé comme incertain.
\newcommand{\uncertain}[1]{\uline{#1}}

% Ajout de l'éditeur — une lettre manquante, un mot sous-entendu.
\newcommand{\add}[1]{\textcolor{blue!50!black}{[#1]}}

% Biffé par Grothendieck lui-même. Ce qu'il a rayé fait partie du document :
% une rature dit souvent où la pensée a changé de direction.
\newcommand{\struck}[1]{\sout{#1}}

% Note du transcripteur, sur l'état matériel du feuillet.
\newcommand{\note}[1]{\par\smallskip{\small\color{gray!60!black}\textit{[#1]}}\par\smallskip}

% Note marginale de Grothendieck, distincte du corps du texte.
\newcommand{\marginal}[1]{\par\smallskip\noindent\hspace{1em}%
  {\small\color{gray!60!black}#1}\par\smallskip}

% --- Titre ------------------------------------------------------------------

\AtBeginDocument{%
  \begin{center}
    {\large\bfseries Fonds Alexandre Grothendieck --- cote n\textsuperscript{o}~\grothFolder}\\[2pt]
    {\itshape\grothFoldertitle}\\[4pt]
    {\small feuillets \grothFirst--\grothLast\ (lot \grothBatch)}\\[2pt]
    {\footnotesize\color{gray!60!black}Transcription établie d'après le fac-similé numérisé
     par l'Université de Montpellier.\\
     Les lectures incertaines sont soulignées, les illisibles notées [\,\ldots\,],
     les ajouts entre crochets.}
  \end{center}
  \vspace{1em}}

\endinput


\folder{89}
\batch{1}
\pages{1}{20}
\dating{s.d. — le groupe « Géométrie et topologie combinatoire » (69 à 102) est daté 1976-[vers 1986]}
\watermark{Édition de démonstration}
\foldertitle{[Hexagone] combinatoire. [Ensembles 3-3] : notes manuscrites (s.d.)}

\begin{document}
\section{Les ensembles 3-3}

\note{ces mots sont de sa main : « les ensembles 3-3 », au crayon, occupent
seuls le haut du feuillet de couverture (page 1), qui ne porte rien d'autre et
ne reçoit donc pas de marqueur de page. L'inventaire les reprend entre
crochets, « [Ensembles 3-3] » ; « [Hexagone] combinatoire » vient de la
première ligne de la page 2}

\subsection{I}

\page{2}

Soit $H = (S, A, R)$ un hexagone combinatoire\note{un hexagone régulier est
dessiné dans la marge gauche, ses six sommets marqués d'un point, avec en
pointillés ses diagonales et ses deux triangles inscrits}. Il lui correspond

\begin{itemize}
  \item[a)] $\omega_H \in \mathrm{Ens}_2$ l'une des deux orientations,
  \item[b)] un torseur $S$ sous
    \struck{$\mathbf{Z}/6\mathbf{Z} \wedge \omega_H \simeq \mathbf{Z}/3\mathbf{Z} \wedge \omega_H$}
    \[
      \mathbf{Z}/6\mathbf{Z} \wedge \omega_H \simeq
      (\mathbf{Z}/3\mathbf{Z} \wedge \omega_H) \times
      \underbrace{(\mathbf{Z}/2\mathbf{Z} \wedge \omega_H)}_{\mathbf{Z}/2\mathbf{Z}}
    \]
\end{itemize}

ce qui revient \uncertain{au même}\note{les deux mots sont écrits par un signe
« $=$ »}, une paire de deux torseurs :

\marginal{Le groupoïde des hexagones combinatoires \uncertain{est} équivalent
à la catégorie $(\mathrm{Ens}_3) \times (\mathrm{Ens}_2)$}

\begin{itemize}
  \item[b$_1$)] un torseur sous $(\mathbf{Z}/3\mathbf{Z} \wedge \omega_H)$ i.e.
    un \struck{tri} élément
    \[
      \begin{cases}
        \Delta \in \mathrm{Ens}_3 & \text{avec iso} \\
        \omega_\Delta \simeq \omega_H
      \end{cases}
    \]
    \marginal{\textbf{NB} a) et b$_1$) équivaut à : la donnée de
    $\Delta \in \mathrm{Ens}_3$}
  \item[b$_2$)] un ens. $T$ à deux éléments.
\end{itemize}

Alors $\Delta$ s'interprète comme l'un des \struck{deux} trois diagonales de
$H$, et $T$ comme l'un des deux triangles inscrits. On récupère $H$ par un
iso. can.
\[
  S \simeq \Delta \times T
\]
(un \ill{} est connu quand on connaît les triangles inscrits et les diagonales
qu'il détermine, et inversement, \uncertain{ceux-ci pouvant} être donnés
arbitrairement ---)

Notons que cette décomposition implique
\[
  S \simeq \coprod_{i \in T} S_i \qquad \text{avec}\quad
  \begin{array}{l} \operatorname{card} T = 2 \\
  \operatorname{card} S_i = 3 \quad \forall\, i \in T \end{array}
\]

\page{3}

en posant
\[
  S_i = \Delta \times \{i\},
\]
et on a un iso. canonique $\omega_{S_i} \simeq \omega_\Delta$ pour
$\forall\, i \in T$, d'où
\[
  \bigwedge_{i \in T} \omega_{S_i} \simeq \mathbf{1}.
  \quad (\text{torseur } \uncertain{\text{neutre}} \text{ sous } \mathbf{Z}/2\mathbf{Z})
\]

Considérons la structure \struck{donnée} sur un ens $S$ de card $6$ donnée par
une partition
\[
  S = \coprod_{i \in T} S_i
\]
du type $(3,3)$, et un \struck{isom.} \supplied{syst. trans. d'}isom. entre les
$\omega_{S_i}$\note{« syst. trans. d' » est écrit au-dessus du mot biffé}, i.e.
un isom
\[
  (\mathrm{A}) \qquad \bigwedge_{i \in T} \omega_{S_i} \simeq \mathbf{1}.
\]

\struck{Considérons les} Le groupe des automorph. de cette structure est
d'ordre $36$ (car en effet

\begin{tikzcd}
1 \arrow[r] & G' \arrow[r] \arrow[d] & G \arrow[r] & \mathfrak{S}_T \arrow[r] \arrow[d, no head, "\shortparallel" description] & 1 \\
 & \prod_{i \in T} \mathfrak{S}_{S_i} \arrow[d, "\prod_{i \in T} \mathrm{sg}_i"] & & \mathbf{Z}/2\mathbf{Z} & \\
 & \mathbf{Z}/2\mathbf{Z} \arrow[d] & & & \\
 & 1 & & &
\end{tikzcd}

)\note{un court trait vertical est tracé au-dessus de $G'$, sans
terme au bout ; le $\prod$ est écrit par-dessus un premier signe}

\struck{La sous-groupe} Une structure hexagonale subordonnée : il

\page{4}

\uncertain{structure unisignée} équivaut à la donnée de l'\textbf{antipodisme}
correspondant $\sigma = \sigma_H$ : cette structure $H$, qui
\struck{\ill{}} \add{est} un autom. \add{involutif} de la structure
\uncertain{unisignée}, induisant la transposition sur $T$ i.e. échangeant
entre eux les $S_i$ --- il y a exactement \ill{} correspondant \ill{} le choix
d'un $i \in T$ aux trois isom
\[
  S_i \simeq S_j \qquad (j \in T,\ j \neq i)
\]
« \uncertain{inverseurs} d'orientations » i.e. rendant
$\omega_{S_i} \simeq \omega_{S_j}$ l'iso. donné grâce à (A). \struck{Soit}

\note{« involutif » est inséré au-dessus de la ligne ; le mot biffé après
« qui » est surchargé d'un « est »}

\marginal{\struck{en effet, \ill{} il \ill{} exactement 2 \ill{}} il faut
exiger de plus que \ill{} trivialise \ill{} \uncertain{commutant aux} \ill{}
$\omega_{S_i}$ ($i \in T$)) ; il y en a trois \ill{} possibilités \ill{} qui
induisent \ill{} des $\omega$, ils correspondent aux éléments \ill{} des
$\bigwedge_{i \in T} S_i$ (par unicité) \ill{}
$\Gamma\omega = \mathbf{Z}/3\mathbf{Z} \wedge_I \omega$ \ill{}}

\note{cette note marginale est écrite en éventail, en biais à une soixantaine
de degrés, sur toute la moitié supérieure de la marge gauche, et chevauche le
début des lignes du corps ; seuls des fragments s'en lisent. Ses trois
premières lignes sont barrées}

Mais l'ens $G_H$ des autom. de $H = H_\sigma$ est le centralisateur de
$\sigma$ \add{$= Z(\sigma)$} --- $G_H$ (est d'indice $3$ dans $G$,
\uncertain{et en particulier le groupe normalisateur}) il est non distingué,
ses trois conjugués correspondent aux trois conjugués de $\sigma$, qui sont
les trois « antipodismes » de $G$, correspondant aux trois structures
hexagonales subordonnées : la structure « trihexagonale » unisignée.

\note{« $= Z(\sigma)$ » est écrit au-dessus de la ligne, et « et en
particulier le groupe normalisateur » dans l'interligne suivant, rattaché par
un trait au mot « distingué »}

Considérons le $2$-groupe (invariant) $V$\note{un $G'$ est écrit au-dessus du
$V$} de $G$ engendré par l'ens. $I$ des trois anti-

\page{5}

podismes $\sigma_i$, on vérifie que pour
\[
  I = \{i, j, k\}
\]
on a
\[
  \sigma_k = \sigma_i \sigma_j \sigma_i
\]
donc (puisqu'on peut exiger les faire en permutation circulaire), donc
\[
  \sigma_i = \sigma_j \sigma_k \sigma_j = \sigma_j (\sigma_i \sigma_j \sigma_i) \sigma_j
  \qquad \text{i.e.} \qquad
  (\sigma_i \sigma_j)^3 = 1
\]

\struck{On constate que $G'$ est le sous-groupe présenté} $G'$ opérant sur
$I$ donne un iso \struck{par}
\[
  \struck{\sigma_i^2 = 1 \quad (i \in I)}
\]
\[
  G' \xrightarrow{\ \sim\ } \mathfrak{S}_I
\]
où $\sigma_i$ correspond à la transposition de $I$ d'indice $i$ --- $G'$ est
présenté par les gén. $\sigma_i$ \add{$(i \in I)$} et les \add{$3 + 6 = 9$}
relations
\[
  \sigma_i^2 = 1 \qquad \sigma_k = \sigma_i \sigma_j \sigma_i
  \qquad (\text{si } I = \{i, j, k\})
\]
\note{« $G'$ opérant sur » est écrit au-dessus de la ligne biffée ; « $(i \in
I)$ » et « $3 + 6 = 9$ » sont insérés au-dessus de la ligne et rattachés par
des traits}
i.e. encore (pour $i, j \in I$ fixés, $i \neq j$), posant
$\sigma = \sigma_i$, $\sigma' = \sigma_j$)
\[
  \struck{\sigma_i^2 = \sigma_j^2 =} \qquad
  \sigma^2 = \sigma'^2 = (\sigma\sigma')^3 = 1 .
\]

Quel est le groupe quotient $G/G'$ (d'ordre $6$) ? Notons que les
\add{données} \struck{\ill{}} de $(S_i)_{i \in T}$
\note{« données » est écrit au-dessus du mot biffé ; la lecture
$(S_i)_{i \in T}$ du dernier terme, en bout de ligne, n'est pas sûre}

\page{6}

et de $\bigwedge_{i \in T} \omega_{S_i} = \mathbf{1}$ équivaut à la donnée d'un
ens. à deux éléments $\omega$ (l'orientation « commune » des $\omega_{S_i}$
identifiés par le syst. transitif d'iso
$\bigwedge_{i \in T} \omega_{S_i} \simeq \mathbf{1}$) et une famille de
\textbf{torseurs} $S_i$ ($i \in T$) sous
$\mathbf{Z}/3\mathbf{Z} \wedge_{\pm 1} \omega = \Gamma_\omega$.
Considérons alors le \supplied{$\Gamma_\omega$-}torseur :
\[
  \mathsf{T} = \struck{\prod_{i \in T}} \bigwedge_{i \in T} S_i
\]
alors $G$ opère sur $(\Gamma_\omega, \mathsf{T})$, je dis que $G'$ y opère
trivialement. En effet, un antipodisme $\sigma$ induisant
$S_i \overset{u}{\simeq} S_j$ et $S_j \overset{u^{-1}}{\simeq} S_i$ induit
\[
  S_i \wedge S_j \xrightarrow{\ \sim\ } S_j \wedge S_i
  \overset{\text{sym}}{\simeq} S_i \wedge S_j
  \qquad (x \wedge y \mapsto (u^{-1} y \wedge u x))
\]
\struck{On a $ux$} Fixons $a \in S_i$, posons $ua = b \in S_j$, on a
\[
  x = \gamma a, \quad y = \gamma' b
\]
donc pour $\gamma, \gamma' \in \Gamma_\omega$
\[
  \sigma(\underbrace{\gamma a \wedge \gamma' b}_{\gamma\gamma'\, a \wedge b})
  = (\underbrace{u^{-1} \gamma' b}_{\gamma' a}) \wedge (\underbrace{u \gamma a}_{\gamma b})
  = \struck{\gamma\gamma'}\ \gamma' a \wedge \gamma b
  = \underset{}{\gamma'\gamma\, a \wedge b}
\]
\note{le $\sigma$ en tête de la formule est écrit en surcharge d'un autre
signe}

\struck{Ainsi} Ainsi $G/G'$ opère dans \struck{\ill{}} $(\Gamma_\omega,
\mathsf{T})$. Mais $\mathsf{T}$ se réalise \struck{comme} \add{canonique} l'un
des \add{$3$} \uncertain{automorph.} involutifs \struck{$\alpha_\psi$} de
$S = \coprod_{i \in T} S_i$ tels que $\alpha$ induise la transposition de $T$
\textbf{et} de $\omega$, et on vérifie que \uncertain{$\alpha_t$ conjugue
$\alpha_{j'}$ en $\alpha_{j''}$} (en fait \ill{} l'a déjà fait \ldots)
\note{« canonique » est écrit au-dessus de « comme », biffé ; le « $3$ » est
inséré au-dessus de la ligne}

\page{7}

\note{la page s'ouvre sur un double trait horizontal, de sa main, au milieu
de la ligne de tête}

Soit $S$ \struck{de} cardinal $6$ muni d'une partition\note{un $R$ est écrit
au-dessus de « partition »} du type $(3,3)$, $T$ l'ens de ses deux classes,
$S_i$ ($i \in T$) ses deux classes. Soit $\mathcal{G}$ le groupe des automorph.
de $S$, donc

\begin{tikzcd}
1 \arrow[r] & \prod_{i \in T} \mathfrak{S}_{S_i} \arrow[r] & \mathcal{G} \arrow[r] & \mathfrak{S}_T \arrow[r] \arrow[d, no head, "\shortparallel" description] & 1 \\
 & & & \mathbf{Z}/2\mathbf{Z} &
\end{tikzcd}

$\mathcal{G}$ d'ordre $72$. Soit $S' \subset \mathcal{G}$ l'ens des éléments
\textbf{involutifs} de $\mathcal{G}$ au-dessus de l'élément $\neq 1$ de
$\mathfrak{S}_T$, i.e. l'ens des scindages de l'ext. précédente.

Pour $i \in T$ fixé, il est en corr. 1-1 avec
\[
  \mathrm{Isom}(S_i, S_j) \qquad \text{où } j \in T,\ j \neq i,
\]
ou encore il est en corr. 1-1 avec l'ens des corr. bijectives entre
$(S_i)_{i \in T}$ i.e. des $\Gamma \subset \prod_{i \in T} S_i$ telles que
$\forall\, i \in T$, $\Gamma \to S_i$ est bij.

\page{8}

C'est un ens qui a $6$ éléments, et on a une \struck{bijection} surjection
\[
  S' \longrightarrow T' \overset{\text{déf}}{=} \bigwedge_i \omega_{S_i}
\]
(notation précédente) qui est de degré $3$. Donc $S'$ est également un ens.
muni d'une partition du type $(3,3)$.

La donnée d'une structure hexagonale équivaut à : celle d'une structure
$(S, R)$, et d'un $s \in S'$ --- celle d'un hexagone et d'un sommet dessus,
\uncertain{à} un

\note{la page s'arrête là, au premier tiers du feuillet ; le reste est blanc}

\subsection{II}

\page{9}

Soit $V$ un \struck{\ill{}} plan vectoriel sur $\mathbf{F}_3$. La droite
projective $P(V)$ a quatre éléments i.e. $V$ \add{contient} quatre droites
vectorielles --- donnons nous une paire $T$ de droites de $V$
$(D_i)_{i \in T}$. On a donc
\[
  (1) \qquad V = \bigoplus_{i \in T} D_i \simeq \prod_{i \in T} D_i
\]
\note{« contient » et le $T$ sont écrits au-dessus de la ligne}

Considérons la paire complémentaire $T' = P(V) \smallsetminus T$,
\struck{\ill{}} on a donc aussi
\[
  (2) \qquad V = \bigoplus_{i \in T'} D_i \simeq \prod_{i \in T'} D_i
\]

Notons les relations
\[
  (3) \qquad
  \begin{cases}
    T' \simeq \bigwedge_{i \in T} D_i^{*} \\
    T \simeq \bigwedge_{i \in T'} D_i^{*}
  \end{cases}
\]
l'application
\[
  \prod_{i \in T} D_i^{*} \longrightarrow T'
\]
qui définit $\bigwedge_{i \in T} D_i^{*} \to T'$ étant
\[
  (4) \qquad (u_i)_{i \in T} \longmapsto \mathbf{F}_3 \cdot \Bigl(\sum_{i \in T} u_i\Bigr).
\]

Si $P$ est un $V$-torseur, i.e. un \uncertain{plan} affine sur $\mathbf{F}_3$
admettant $V$ comme groupe de translations, donc les formules (1) et (2)
définissent respectivement deux décompositions de $P$ en produits
\[
  (5) \qquad P \simeq \prod_{i \in T} P_i\, , \qquad P \simeq \prod_{i \in T'} P_i
\]
\note{le mot rendu par « plan » est barré d'un trait, et rien ne le remplace}

\page{10}

\struck{où pour} pour tout $i \in T \amalg T' = P(V)$, $P_i$ est un
$D_i$-torseur, défini par
\[
  (6) \qquad P_i = P / D_{\bar\imath} = P \wedge_V (D_i \simeq V / D_{\bar\imath})
\]
où $i \mapsto \bar\imath$ est l'\struck{élément} \add{involution} de $P(V)$ qui
\struck{est} correspond à la partition du type $(2,2)$ $\{T, T'\}$, qui dans $T$
et dans $T'$ se réduit à la transposition dans ces ensembles. Ainsi,
l'ensemble $P$, de cardinal $9$, est muni, via
\[
  (\ast) \qquad P \simeq \prod_{i \in T} P_i
\]
d'\add{une paire de} \struck{deux} partitions du type $(3,3,3)$, indexées par
$T$, ces partitions étant « orthogonales » en ce sens que l'application
canonique \add{de} $P$ \uncertain{vers} le produit des deux ens. quotients
correspondants est bijective, ce qui signifie aussi que pour toute classe de
l'une et toute classe de l'autre, l'intersection a exactement un élément (il
suffit de supposer qu'elle est toujours non vide, ou qu'elle a toujours au
plus un élément \ldots). Et de même ensemble $P$ est muni d'une deuxième paire
(distincte de la première) de partitions de type $(3,3)$\note{sic : le
feuillet porte $(3,3)$ et non $(3,3,3)$}, « orthogonales », indexée par $T'$,
correspondant à la décomposition
\[
  (\ast') \qquad P \simeq \prod_{i \in T'} P_i
\]
\note{« involution », « une paire de », « de » et « orthogonales », indexée
par $T'$ » sont écrits au-dessus des lignes}

\page{11}

Ceci étant, on a des équivalences de catégories

\begin{tikzcd}[column sep=small]
 & (\mathbf{Tors}(V)) & \\
\prod_{i \in T} (\mathbf{Tors}(D_i)) \arrow[ur, leftrightarrow, "\approx"] & & \prod_{i \in T'} (\mathbf{Tors}(D_i)) \arrow[ul, leftrightarrow, "\approx"']
\end{tikzcd}

(7)\note{le numéro (7) est porté à gauche du diagramme ; son
soulignement de « Tors » est rendu en gras}

Notons que pour une droite vectorielle $D$ sur $\mathbf{F}_3$, la catégorie
$\mathbf{Tors}(D)$ est isomorphe canoniquement à celle des ens. $E$ de cardinal
$3$, munis d'une bijection
\[
  \omega_E \simeq D^{*}
\]
\marginal{$\omega_E$ : l'une des $2$ orientations de $E$, i.e. des deux
permutations circulaires de $E$ inverses l'une de l'autre}
de sorte que la catégorie $\prod_{i \in T} (\mathbf{Tors}(D_i))$ peut être
regardée comme celle des familles à deux éléments, indexées par $T$,
d'ensembles de cardinal $3$
\[
  (P_i)_{i \in T}
\]
munis chacun d'une bijection
\[
  \omega_{P_i} \simeq D_i^{*} \qquad (i \in T)
\]
--- et de même pour la catégorie $\prod_{i \in T'} (\mathbf{Tors}(D_i))$.
\note{la glose sur $\omega_E$ est écrite en petits caractères sous la formule,
et non en marge}

\page{12}

On peut aussi l'interpréter, \add{via
$(P_i)_{i \in T} \longmapsto \coprod_{i \in T} P_i$,} comme la catégorie des
ensembles $E$ de cardinal $6$, munis d'une partition \add{$(P_i)_{i \in T}$}
du type $(3,3)$ indexée par $T$, et pour chaque $i \in T$ d'une bijection
$\omega_{P_i} \simeq D_i^{*}$. Ainsi les deux équivalences (7), fournissent une
équivalence
\note{« via $(P_i)_{i \in T} \mapsto \coprod_{i \in T} P_i$ » et
« $(P_i)_{i \in T}$ » sont écrits au-dessus des lignes}
\[
  (8) \qquad
  \begin{array}{l}
    \text{Catégorie des ens. } E \text{ de} \\
    \text{cardinal } 6, \text{ munis d'une} \\
    \text{partition } (P_i)_{i \in T} \text{ du type } (3,3), \\
    \text{indexée par } T, \text{ et d'une} \\
    \struck{\text{famille}}\ \text{paire de bijections} \\
    \qquad D_i^{*} \simeq \omega_{P_i}
  \end{array}
  \quad \xrightarrow{\ \approx\ } \quad
  \begin{array}{l}
    \text{Catégorie analogue} \\
    \text{définie en termes} \\
    \text{de } T' \text{ et de} \\
    \quad (D_i^{*})_{i \in T'}
  \end{array}
\]

Quand \add{(\uncertain{grâce à} $V$)} on échange le rôle de $T$ et de $T'$, il
est évident que le diagramme d'équivalences (7) reste le même à symétrie près,
et que l'équivalence (8) est donc remplacée par \struck{l'} une équivalence,
\add{conséquemment} quasi-inverse de la précédente.

Les catégories figurant de part et d'autre dans (8), se définissent en termes
de la famille, indexée par $P(V) = T \amalg T'$, d'ensembles à deux éléments
$\omega_i = D_i^{*}$ (dont la donnée équivaut à celle de la famille des droites
vect. $D_i$ \ill{})

\page{13}

vu que \struck{$\mathbf{F}_3^{*}$} l'on a
\[
  \mathbf{F}_3^{*} = \{\pm 1\} \quad ).
\]

\struck{Ainsi} \add{Notons que} la donnée de $\{(\omega_i)_{i \in T}\}$ permet
ainsi de reconstituer \add{(à isom. près)} la famille $(D_i)_{i \in T}$ de
droites vectorielles sur $\mathbf{F}_3$ (via $D_i = \omega_i \amalg \{0\}$
\ldots), donc le vectoriel
\[
  V = \prod_{i \in T} D_i
\]
sur $\mathbf{F}_3$, donc aussi \add{le plan vectoriel}
$V = \prod_{i \in T} D_i$ avec la paire de droites indexée par $T$, d'où la
paire complémentaire $(D_i)_{i \in T'}$. \struck{La}

\struck{Donc ainsi on fait} donnée d'une famille $(\omega_i)_{i \in T}$ d'ens.
à deux éléments, indexée par $T$, équivaut à celle d'un \textbf{carré}
\add{comb.} $Q$, admettant $T$ comme ens. des diagonales. On trouve
\add{ainsi} \struck{une} équivalences
\[
  (9) \qquad
  \begin{array}{c} \text{Carrés combinatoires} \\ Q \end{array}
  \quad \overset{\approx}{\longleftrightarrow} \quad
  \begin{array}{l}
    \text{plans vectoriels } V \text{ sur } \mathbf{F}_3, \\
    \text{munis d'une paire} \\
    \text{de droites vectorielles} \\
    \quad (D_i)_{i \in T}
  \end{array}
\]
\note{« Notons que », « (à isom. près) », « le plan vectoriel », « comb. »,
« ainsi » et le $V$ de (9) sont écrits au-dessus des lignes. La lecture
« carré » d'un mot souligné qu'on pourrait prendre pour « couvé » est dictée
par la suite : un ensemble de quatre sommets dont $T$ est l'ensemble des deux
diagonales}

dans cette équivalence, donc, \struck{on $T$ correspond à}
\[
  (10) \qquad
  \left\{
  \begin{array}{l}
    T \simeq \operatorname{diag} Q
    \quad \text{et} \quad D_i^{*} \simeq d_i \quad \text{pour } i \in T \\[4pt]
    S_Q = \bigcup_{i \in T} D_i^{*} = \coprod_{i \in T} D_i^{*} \\[4pt]
    A_Q = \prod_{i \in T} D_i^{*} \simeq V^{*} \smallsetminus \bigcup_{i \in T} D_i^{*}
    = \bigcup_{i \in T'} D_i^{*}
  \end{array}
  \right.
\]
où $d_i$ est l'ensemble des deux sommets de $Q$ qui constituent la diagonale
$i$
\note{à gauche de l'accolade, un « Diag $Q$ » est biffé au-dessus du numéro
(10). Le dernier membre « $= \bigcup_{i \in T'} D_i^{*}$ » est rejeté dans la
marge gauche, en bas du feuillet, qui coupe la dernière ligne ; dans la
ligne de $A_Q$ un premier signe est surchargé en $\prod$}

\page{14}

\struck{Rien} On a d'ailleurs
\[
  \underset{\substack{\text{ens. d'arêtes} \\ \text{de } Q}}{A_Q}
  \simeq
  \underset{\substack{\text{produit} \\ \text{des deux} \\ \text{diagonales de } Q}}{\prod_{i \in T} D_i^{*}}
  \simeq V^{*} - \bigcup_{i \in T} D_i^{*} = \bigcup_{i \in T'} D_i^{*}
\]
ce qui montre que \struck{\ill{}} le carré défini par l'\struck{ens.}
\add{paire} $(D_i)_{i \in T'}$ des droites de $V^{*}$ complémentaire de la paire
$(D_i)_{i \in T}$ s'identifie au carré dual $DQ$ de $Q$. Si on désigne par
$V(Q)$ le \add{plan} vectoriel sur $\mathbf{F}_3$ défini par $Q$, on trouve
donc un isomorphisme
\[
  (11) \qquad V(DQ) \xrightarrow[\ \sim\ ]{\ \alpha_Q\ } V(Q)
\]
et de façon symétrique
\[
  (\uncertain{12}) \qquad
  \underset{\substack{\wr \\ DDQ}}{V(Q)} \xrightarrow{\ \alpha_{DQ}\ } V(DQ)
\]
\note{le numéro de cette formule est écrit comme celui de la précédente ; la
suite, (13), demande (12). Le $Q$ de $V(Q)$ est surchargé}

Ceci posé, on a la relation
\[
  (13) \qquad \alpha_Q\, \alpha_{DQ} = -\mathrm{id}_{V(Q)}
  = V(\underset{\substack{\text{antipodisme} \\ \text{de } Q}}{a_Q})
\]
\note{les indices du premier membre sont écrits en surcharge d'une première
version}
et de même bien entendu
\[
  \alpha_{DQ}\, \alpha_Q = -\mathrm{id}_{V(DQ)} = V(a_{DQ})
\]

\page{15}

Soient $V$ un Module \struck{loc.} libre \add{de rang $2$} sur un
\add{anneau $\Lambda$} \struck{topos loc. annelé}, d'où $P(V) =$ ens. des
« droites vectorielles » de $V$ (i.e. des sous-modules \struck{libres de}
\add{loc.} inversibles facteurs directs), \struck{et $P(V)^! =$}
\begin{itemize}
  \item[(1)] $V^{*} =$ ens. des éléments de $V$ partout $\neq 0$ (i.e. qui font
    partie d'une base)
  \item[(2)] $P(V)^! = V^{*} / \Lambda^{*} \hookrightarrow P(V)$
    l'ens. des « droites dans $V$ » qui sont des modules libres de rang $1$.
\end{itemize}
\note{« de rang $2$ », « anneau $\Lambda$ » et « loc. » sont écrits au-dessus
des mots biffés}

On a envie d'expliciter comment on récupère \add{le $\Lambda$-module} $V$ à
l'aide de l'ens. $P(V)^!$, muni d'une structure supplémentaire, qui est
\struck{un} essentiellement le $\Lambda^{*}$-\uncertain{torseur} relatif
(ensembliste) $V^{*}$ sur $P(V)^! = X$.

Pour $x \in P(V)^!$, soit $D_x \subset V$ la droite vectorielle correspondante,
et posons
\[
  (3) \qquad T_x = D_x^{*} \qquad (\Lambda^{*}\text{-torseur})
\]
Soit de plus
\[
  (4) \qquad \omega_V = \det(V)^{*} \qquad (\text{un } \Lambda^{*}\text{-torseur})
\]

Sur $P(V)^!$ on a la relation binaire de disjonction
\[
  x, y \in P(V)^! \ \text{« \textbf{disjoints} »} \ \text{ssi} \ V \simeq D_x \oplus D_y
\]
(si $\Lambda$ est un corps, cela équivaut à : $x \neq y$ ; si $\Lambda$ est
local, cela équivaut à : $x_0 \neq y_0$ (l'indice $0$

\page{16}

désignant la réduction mod. l'idéal maximal). Si $x$ et $y$ sont disjoints,
alors \add{(et alors seulement)} \struck{l'appl.}
\[
  \begin{array}{c}
    (u, v) \longmapsto u \wedge v \\
    D_x \times D_y \longrightarrow \det V
  \end{array}
\]
induit
\[
  T_x \times T_y \longrightarrow \omega_V
\]
qui se factorise par un \supplied{iso}morphisme
\[
  (5) \qquad
  \underset{\substack{\text{produit contracté} \\ \text{de } \Lambda^{*}\text{-torseurs}}}{T_x \wedge T_y}
  \xrightarrow{\ \sim\ } \omega_V
\]

\struck{noté} L'image de $(u \wedge v)$ \add{dans $\omega_V$} par cette
application sera encore notée $u \wedge v$. On a encore
\[
  (6) \qquad u \wedge v = - v \wedge u
\]
\marginal{\textbf{NB} $-1 \in \Lambda^{*}$, et on pose, pour un
$\Lambda^{*}$-torseur $T$, et $u \in T$,
$-u \overset{\text{déf}}{=} (-1).u$)}

Soient $x, y$ disjoints \add{$\in P(V)^!$}, et $z \in P(V)^!$ quelc., soient
\[
  u \in T_x, \quad v \in T_y, \quad w \in T_z
\]
considérons les relations
\[
  (7) \qquad u \wedge v = v \wedge w = w \wedge u
\]
équivalentes à
\[
  (7') \qquad w \wedge v = v \wedge u = u \wedge w \qquad )
\]
\note{le premier membre de (7') est écrit en surcharge}
on vérifie de suite qu'elle équivaut à
\[
  (8) \qquad u + v + w = 0
\]
et qu'elle implique que $z$ est disjoint de $x, y$.

\page{17}

Je dis que \add{le $\Lambda$-module} $V$ se reconstitue \add{(du moins si
$\Lambda$ est un corps)} à partir de
\begin{itemize}
  \item[a)] l'ens. $P = P(V)^!$
  \item[b)] la famille des \add{$\Lambda^{*}$-torseurs} $(T_x)_{x \in P}$ ---
    i.e. le \supplied{$\Lambda^{*}_P$-}torseur relatif
    $\coprod_{x \in P(V)^!} T_x = V^{*}$ sur $P$
  \item[c)] le $\Lambda^{*}$-torseur $\omega = \omega_V$
  \item[d)] la donnée de la relation de disjonction dans $P = P(V)^!$
    [ou ce qui revient au même, dans $V^{*}$ --- cette relation étant
    \supplied{bi-}stable par l'opération de $\Lambda^{*}$ (i.e. $u, v$ disjoints
    $\Rightarrow$ $\lambda u, \mu v$ disjoints pour $\lambda, \mu \in
    \Lambda^{*}$)].
  \item[\struck{e)}] \struck{le $\Lambda^{*}$-torseur $\omega = \omega_V$}
  \item[e)] la donnée, pour $x, y$ disjoints, de l'isom. de
    $\Lambda^{*}$-torseurs
    \[
      T_x \wedge T_y \simeq \omega \qquad \text{noté}\quad
      u \mathbin{\underset{\Lambda^{*}}{\wedge}} v \longmapsto u \wedge v .
    \]
\end{itemize}
\note{à gauche de a) et b), qu'embrasse un trait vertical, écrit en biais
dans la marge :}
\marginal{équivaut à : la donnée de l'ens. $V^{*}$, avec l'opération (libre)
de $\Lambda^{*}$ sur $V^{*}$}
\note{à gauche de d) et e), qu'embrasse un crochet, de même :}
\marginal{Si $\Lambda$ est un corps, cette relation est la relation triviale
$x \neq y$, et on \ill{} de donnée \ill{}}

En effet, pour \struck{\ill{}} $u, v, w$ dans $V^{*}$, deux à deux disjoints,
on trouve
\[
  u + v + w = 0 \iff u \wedge v = v \wedge w = w \wedge u
\]
donc on arrive à trouver la loi de composition \add{partielle}
\[
  (u, v) \longmapsto u + v \quad \text{sur } V^{*} \text{ pour } u, v
  \ \textbf{disjoints}.
\]

\struck{On a} Dans le cas où $\Lambda$ est un corps, posant
\[
  \begin{cases}
    V = V^{*} \amalg \{0\} \\
    D_x = T_x \cup \{0\} \subset V
  \end{cases}
\]
les $D_x$ sont des droites vectorielles sur $\Lambda$, et on

\page{18}

définit $u + v$ ($u, v \in V^{*}$) séparément quand $u$ ou $v = 0$ (comme
étant $v$ resp. $u$), quand $u, v \in V^{*}$ et $u, v$ dans une même $D_x$
(grâce à l'addition dans $D_x$), et quand $u, v$ sont disjoints par la
\struck{\ill{}} somme $-w$, où $w$ est l'\textbf{unique} $w \in V^{*}$ qui
soit disj. de $u, v$ et tel que
\[
  u \wedge v = v \wedge w = w \wedge u .
\]
Pour pouvoir réaliser cette construction, on a besoin des \struck{ceci}
propriétés que voici sur les données a) b) c) [d)] e) :
\begin{itemize}
  \item[a)] Pour $x, y, z \in P$ mutuellement disjoints, on peut trouver
    $u \in T_x$, $v \in T_y$, $w \in T_z$ tels que l'on ait
    \[
      u \wedge v = v \wedge w = w \wedge u \qquad (\ill{})
    \]
    (et alors on peut même supposer que $u$ ou $v$ ou $w$ \struck{\ill{}} soit
    donné d'avance).
\end{itemize}
\note{« $\in P$ » et les deux « ou » sont écrits au-dessus des lignes ; la
parenthèse au bout de la formule reste illisible. La page s'arrête au tiers
inférieur du feuillet}

\subsection{Polygones}

\page{19}

$\mathrm{Pol}\ \mathrm{top} \longrightarrow \mathrm{Pol}\ \mathrm{comb}$
\add{est} ess.\ surjectif\supplied{, \uncertain{\textbf{Foncteur réalisation topologique}}} : dim. \ill{}
\note{à droite, un polygone dessiné : un cercle portant alternativement des
points et des croix, un arc épaissi entre deux points}

dualité : échange de $S$ et $A$

point de vue ens. ordonnés $P = S \amalg A$
\struck{\ill{}} ($S$ él. minimaux, $A$ éléments maximaux
\begin{itemize}
  \item[] chaque \struck{élément} \uncertain{minimal} majoré \ill{}
  \item[] chaque élément minimal majoré \add{\uncertain{strict.}} par
    exactement deux él. maximaux)
  \item[] chaque élément maximal minoré par exactement deux éléments
    \uncertain{minimaux})
\end{itemize}
ici dualité : renversement de la relation d'ordre

\struck{Point de vue} \add{Composantes connexes des \uncertain{v.t.} des ens.
comb. polygones combinatoires (définitions combinatoires)} Point de vue
\textbf{ens. de base $S$} (valable si les comp. connexes sont d'ordre
$\geqslant 3$)
\[
  A \hookrightarrow \mathfrak{P}_2(S) \quad
  \text{structure ens. de parties à deux éléments de } S
  \ (\text{appelées \textbf{arêtes}})
\]
tel que $\forall\, s \in S$ \uncertain{appartient} : exactement deux arêtes

Point de vue dual : structure sur $A$, d'une ens.
$S \hookrightarrow \mathfrak{P}_2(A)$ (qui \ill{} « les
\uncertain{sommets} »)

Point de vue ordres circulaires (cas des polygones d'ordre $\geqslant 3$) :
donnée sur $S$ (ou $A$) d'une paire d'ordres circulaires \add{(i.e.
permutations transitives)} inverses l'un de l'autre --- une
\textbf{orientation} de $\Pi$ correspond au choix de l'un des deux ordres
circulaires en question.

\note{les deux lignes insérées au-dessus de « Point de vue ens. de base
$S$ » sont serrées et en partie illisibles}

\page{20}

Orientation d'une \textbf{arête} : par choix d'une « origine ».
Orientations induites sur le bord de l'arête (extrémité $-$ origine)

Orientation d'un contour combinatoire : choix d'une orientation sur chaque
arête, de telle façon que pour tout sommet $s$, les deux orientations induites
sur $s$ soient opposées \quad $\longrightarrow\!\bullet\!\longrightarrow$
\marginal{orientations opposées}

\struck{\ill{}} Relations avec le cas topologique ; orientation d'un segment
topologique, d'une $1$-variété topologique

\textbf{Théorème} a) Tout \struck{\ill{}} polygone combinatoire admet
exactement $2$ orientations, \struck{\ill{}} \add{\uncertain{les} \ill{}
d'un contour.}
\marginal{i.e. il y a \ill{} orientations}

b) À toute orientation $\omega$ \add{d'un contour,} correspond une unique
\uncertain{bijection} $f_\omega$ telle que
\begin{itemize}
  \item[1)] $f_\omega$ conserve l'orientation
  \item[2)] $\forall\, s \in S$, $\struck{\ill{}}\,(s, f_\omega(s))$ soit une
    \textbf{arête} $a$, \add{et} dont l'origine $o_\omega(a)$ soit $s$.
  \item[\struck{c)}] \struck{\ill{}}
\end{itemize}
\textbf{NB} 2) caractérise déjà $f_\omega$, si les comp. de $\Pi$ sont d'ordre
$\geqslant 3$.

c) $\omega \mapsto f_\omega$ induit une bij. entre l'ens. des orientations de
$\Pi$, et l'ens. des \uncertain{automorphismes} de $\Pi$ qui satisfont 2).

\note{la marge gauche porte, en biais et reliée au texte par un long trait
coudé, une seconde formulation de 2) :}
\marginal{b) $\forall\, s \in S$, $\ill{}\, f(s)$ soit adjacent à $s$ ;
$\beta$) $\forall\, a \in A$, $\ill{}$ ; $f(a) \neq a$ \ill{} origine \ill{}
de $f(s)$ \ill{}}

\textbf{NB} La description des orientations, via $f$ est \ill{}. Orientons
$\Pi$ en un sommet $s$ revient à choisir l'arête « origine » et l'arête
« extrémité » en $s$ \quad (croquis : deux segments se rencontrant en un
point, marqués $+$ et $-$)
\note{les trois lignes biffées qui précèdent le premier \textbf{NB} ne se
lisent pas ; le feuillet s'arrête en bas de page, la suite manque}

\end{document}
